\section{制度的连续，不等于政治没有转向}

1048年李元昊在内政冲突中遇害，西夏转入幼主与后族主政，提醒宋廷西北面对的并非静止的敌国；同年黄河改道所造成的河北压力，也让边防、河工和财政彼此纠缠。1057年嘉祐科举使苏轼、苏辙等进入中央政治视野，科举榜单可以确认，然而它不能说明地方士人已共享一种文风或政治立场。考试、学校、荐举和家族网络共同构成官僚再生产，任何一张榜单都只是其中一个入口。

1063年仁宗去世，1064年濮议展开，1067年神宗继位，1068年王安石进入中枢；1069年新法陆续推行，1071年免役法在部分地区施行。关于继承、任命与法令，《宋史》《续资治通鉴长编》和《宋会要辑稿》能够互相校订；但它们回答的层次不同：本纪建立朝廷年序，长编保存密集的政事线索，会要记录制度条目。三者都不能直接把颁行等同于所有州县的实际执行。

1074年灾荒与政论交织，反对者以民间困乏批评新法；1075年宋军在熙河经营堡寨、道路与羁縻关系；1079年苏轼因诗案受审。灾异、军事与文字案件进入同一政治时期，却不构成一条单因果链。灾异语言可以批评政府，堡寨要求人力与粮道，诗文案则界定表达边界；它们分别说明改革如何在乡村、边境与官僚文化中遇到不同的摩擦。

\section{党籍、征发与危机前的国家能力}

1080年永乐城陷落，1081年多路伐夏未获预期，1084年《资治通鉴》进呈，1085年神宗去世，1086年元祐政府废改新法，1093年哲宗亲政，1094年绍圣政府恢复、调整新法。战争失利与政策更替反复提醒朝廷：军费、基层征发和人事联盟都无法由一项法令永久解决。史书编纂在这里也不是退离政治的文化余事；它提供了论辩治道和评价前代的共同语言，但并不裁定当时政策孰是孰非。

1097年苏轼被贬儋州，1100年端王继位为徽宗，1117年两浙征发与花石纲问题成为方腊起事的背景之一。贬谪地理、宫廷继承和东南征发看似彼此遥远，却共同揭示中央命令的不同落点。关于花石纲与起事，镇压方记录尤其容易把参与者的口号、宗教和人数写成可控的类别；较可靠的是，徭役、运输、乡里网络和东南税源的紧张确实彼此交织。

1119年宋金继续讨论夹攻辽与燕云，1124年张觉事件未能消除互疑，1125年金军南下时宋的堡寨、勤王与信息体系相继承压。1126年康王赵构在围汴和勤王活动中仍保有行动空间，为其后南宋朝廷重建留下条件。把靖康只解释为某位君主或某一条盟约的失败，会遮蔽这一系列累积：边防的情报与供给、政治分类造成的协作困难，以及东南资源被如何征发和转运，都限制了危机中的选择。
